%Chargement des paquets
\usepackage{amsmath}
\usepackage{amsfonts}
\usepackage{amssymb}
\usepackage{enumerate}
\usepackage{mathtools}
\usepackage{bbm}
\usepackage{xparse, etoolbox}
\usepackage{enumerate}
\usepackage{enumitem}
\usepackage{mathabx}
\usepackage{babel}[french]

%environment
\usepackage{amsthm}
\usepackage{keytheorems}
\usepackage{hyperref} 

%Interface théorème
\newkeytheorem{lem}[
	name = Lemme
]

\newkeytheorem{prop}[
	name = Proposition
]

\newkeytheorem{thm}[
	name = Théorème
]

\newkeytheorem{coro}[
	name = Corollaire
]

\renewcommand*{\proofname}{Démonstration}

\theoremstyle{definition}
\newtheorem{defn}{Définition}

\theoremstyle{definition}
\newtheorem*{nota}{Notation}

\theoremstyle{definition}
\newtheorem*{limi}{Liminaire}

\theoremstyle{remark}
\newtheorem{rmq}{Remarque}

\theoremstyle{plain}
\newtheorem*{fait}{Fait}


%Ensembles de nombres
\newcommand{\N} {\mathbb{N}}
\newcommand{\Ne}{\N^\ast}
\newcommand{\Z} {\mathbb{Z}}
\newcommand{\Q} {\mathbb{Q}}
\newcommand{\R} {\mathbb{R}}
\newcommand{\Rb}{\overline{\mathbb{R}}}
\newcommand{\Rp}{\R_+}
\newcommand{\Rm}{\R_-}
\newcommand{\K} {\mathbb{K}}
\newcommand{\Ket} {\mathbb{K^{\ast}}}
\newcommand{\Cx}{\mathbb{C}}

%Ensembles, tribus
\newcommand{\A}{\mathbb{A}}
\renewcommand{\P}{\mathcal{P}}
\newcommand{\T}{\mathcal{T}}
\newcommand{\F}{\mathcal{F}}
\newcommand{\C} {\mathcal{C}}
\newcommand{\ouv}{\mathcal{O}}
\newcommand{\B}{\mathcal{B}}
\newcommand{\M}{\mathcal{M}}
\newcommand{\etag}{\mathcal{E}}
\newcommand{\etagOp}{\etag^{+}(\Omega)}
\newcommand{\etagO}{\etag(\Omega)}
%%Lp avant quotientage
\newcommand{\Lic}{\mathcal{L}^1}
\newcommand{\Lpc}{\mathcal{L}^p}
\newcommand{\Linfc}{\mathcal{L}^{\infty}}
\newcommand{\Lifull}{\Lic(\Omega, \B, m)}
\newcommand{\Lpfull}{\Lpc(\Omega, \B, m)}
\newcommand{\Linffull}{\Linfc(\Omega, \B, m)}
%%Lp après quotientage
\newcommand{\Li}{L^1}
\newcommand{\Ld}{L^2}
\newcommand{\Lp}{L^p}
\newcommand{\Linf}{L^{\infty}}
\newcommand{\Listd}{\Li(\Omega, m)} 
\newcommand{\Ldstd}{\Ld(\Omega, m)} 
\newcommand{\Lpstd}{\Lp(\Omega, m)} 

%Quelques opérateurs
\DeclareMathOperator{\codim}{codim}
\DeclareMathOperator{\rg}{rg}
\DeclareMathOperator{\tr}{tr}
\DeclareMathOperator{\vect}{Vect}
\DeclareMathOperator{\ima}{Im}
\DeclareMathOperator{\id}{Id}
\DeclareMathOperator{\sign}{sign}
\DeclareMathOperator{\ext}{ext}
\newcommand{\ind}{\mathbbm{1}}
\newcommand*{\spr}[2]{\left(\,#1\,\middle|\,#2\,\right)}
\newcommand{\interior}[1]{%
  {\kern0pt#1}^{\mathrm{o}}%
}
\DeclareMathOperator{\card}{card}
\DeclareMathOperator{\ppcm}{ppcm}

%Commandes de pure flemme
\newcommand{\veps}{\varepsilon}
\newcommand{\intab}{\int_{a}^{b}}
\newcommand{\intR}{\int_{-\infty}^{\infty}}
\newcommand{\intinf}{\int_{- \infty}^{+ \infty}}
\newcommand{\equi}{\Leftrightarrow}
\newcommand{\Kev}{$\K$-espace vectoriel }
\newcommand{\Kevs}{$\K$-espaces vectoriels }
\newcommand{\Rev}{$\R$-espace vectoriel }
\newcommand{\Revs}{$\R$-espaces vectoriels }
\newcommand{\Cev}{$\Cx$-espace vectoriel }
\newcommand{\Cevs}{$\Cx$-espaces vectoriels }
\newcommand{\evn}{(E, \norm{.})}
\newcommand{\interf}[2]{[#1,#2]}
\newcommand{\limb}{\lim\limits_}
\newcommand{\lebn}{\lambda_n}
\newcommand{\pp}{\text{~presque partout}}
\newcommand{\derivp}[2]{\frac{\partial #1}{\partial #2}}
\newcommand{\famiu}{(u_i)_{i \in I}}
\newcommand{\sev}{\text{~s.e.v.}}
\newcommand{\Mnp}{M_{n,p}(\K)}
\newcommand{\Mn}{M_{n}(\K)}
\newcommand{\tq}{\text{~t.q.~}}
\newcommand{\mq}{~montrer~que~}
\newcommand{\et}{\quad \text{et} \quad}
\newcommand{\ou}{\text{~ou~}}
\newcommand{\Kx}{\K[X]}


%Commandes de gars chelou
\renewcommand{\subset}{\subseteq}

%Commande restriction, ty duckduckgo
\def\restr#1#2{\mathchoice
              {\setbox1\hbox{${\displaystyle #1}_{\scriptstyle #2}$}
              \restraux{#1}{#2}}
              {\setbox1\hbox{${\textstyle #1}_{\scriptstyle #2}$}
              \restraux{#1}{#2}}
              {\setbox1\hbox{${\scriptstyle #1}_{\scriptscriptstyle #2}$}
              \restraux{#1}{#2}}
              {\setbox1\hbox{${\scriptscriptstyle #1}_{\scriptscriptstyle #2}$}
              \restraux{#1}{#2}}}
\def\restraux#1#2{{#1\,\smash{\vrule height .8\ht1 depth .85\dp1}}_{\,#2}} 

%Quelques normes
\DeclarePairedDelimiter\abs{\lvert}{\rvert}
\DeclarePairedDelimiter\labs{\bigl\lvert}{\bigr\rvert}
\DeclarePairedDelimiter\norm{\lVert}{\rVert}
\DeclarePairedDelimiterXPP\normi[1]{}\lVert\rVert{_1}{\ifblank{#1}{\:\cdot\:}{#1}}
\DeclarePairedDelimiterXPP\normd[1]{}\lVert\rVert{_2}{\ifblank{#1}{\:\cdot\:}{#1}}
\DeclarePairedDelimiterXPP\normp[1]{}\lVert\rVert{_p}{\ifblank{#1}{\:\cdot\:}{#1}}
\DeclarePairedDelimiterXPP\normq[1]{}\lVert\rVert{_q}{\ifblank{#1}{\:\cdot\:}{#1}}
\DeclarePairedDelimiterXPP\norminf[1]{}\lVert\rVert{_\infty}{\ifblank{#1}{\:\cdot\:}{#1}}

%Bracket en angle
\DeclarePairedDelimiter\angl{\langle}{\rangle}

%Set, parenthèses
\newcommand{\set}[1]{\{ #1 \}}
\newcommand{\bigset}[1]{\bigl\{ #1 \bigr\}}
\newcommand{\Bigset}[1]{\Bigl\{ #1 \Bigr\}}
\newcommand{\bigpar}[1]{\bigl( #1 \bigr)}
\newcommand{\Bigpar}[1]{\Bigl( #1 \Bigr)}

%Opitmisation des opérateurs de comparaison
\DeclarePairedDelimiter\tio{o (}{)}
\DeclarePairedDelimiter\gro{O (}{)}


\pagestyle{fancy}

\fancyhead[C]{Théorème des restes Chinois}
\fancyhead[R]{}

\begin{document}

\underline{Sources :} Rombaldi - Algèbre - page 249. \newline

\underline{Leçons :} 
\begin{itemize}
\item 120 : Anneaux Z/nZ. Applications.
\item 122 : Anneaux principaux. Exemples et applications.
\item 142 : PGCD et PPCM, algorithmes de calcul. Applications.
\end{itemize}
\vspace{10pt}

Ce résultat date du IIIème siècle, il est d'abord formulé par le philosophe chinois Sun Zi (d'où le nom du théorème), 
puis reprit par le mathématicien chinois Qin Jiushao sous cette forme :

\begin{center}
\textit{Soient des objets en nombre inconnu. Si on les range par 3 il en reste 2. Si on les range par 5, il en reste 3 et si on les range par 7, il en reste 2. Combien a-t-on d'objets ?}
\end{center}

On propose une généralisation de ce résultat dans un anneau principal, puis une application dans les anneaux quotients $\Z / n\Z$.

\section{Le théorème}

\begin{nota}
Soit $A$ un anneau principal, et $(a_j)_{1 \leq j \leq r}$ une suite de $r \geq 2$ éléments non nuls et non inversibles de $A$. On note :
\[ a = \prod_{j=1}^r a_j . \]
Pour tout indice $j$, $\pi_j$ désigne la surjection canonique de $A$ sur l'anneau quotient $A / (a_j)$. \\
On muni le produit cartésient $\prod_j A / (a_j)$ d'une structure d'anneau commutatif unitaire en posant :
\[ \left \{ \begin{array}{c c c} 
(\pi_1(x_1), \dots, \pi_r (x_r)) + (\pi_1(y_1), \dots, \pi_r (y_r)) & = & ( \pi_1(x_1+y_1), \dots, \pi_r(x_r+y_r) \\
(\pi_1(x_1), \dots, \pi_r (x_r)) \cdot (\pi_1(y_1), \dots, \pi_r (y_r)) & = & ( \pi_1(x_1\cdot y_1), \dots, \pi_r(x_r \cdot y_r)
\end{array} \right. \]
(autrement dit, on définit l'addtion et la multiplication composante par composante). \\
On note $\pi$ la surjection canonique de $A$ sur l'anneau quotient $A/ (a)$. \\
Enfin, on désigne par $(b_j)_{1 \leq j \leq r}$ la suite d'éléments de $A$ définie par :
\[ b_j = \dfrac{a}{a_j} = \prod_{i \neq j} a_i . \]
\end{nota}

\begin{lem}
Si les éléments $a_j$ sont deux à deux premiers entre eux, alors les $b_i$ sont premiers entre eux dans leur ensemble.
\end{lem}

\begin{proof}
Supposons que les $b_i$ ne soient pas premiers entre eux dans leur ensemble et exhibons $d$, 
élément premier de $A$ et diviseur commun 
des $b_i$ (ceci est possible dans un anneau principal, donc en particulier factoriel). 
En particulier, comme $d$ divise $b_1$, par primalité il divise au moins un $a_j$ pour $j \geq 2$. 
Mais comme $d$ divise aussi $b_j$, il va diviser deux éléments$a_j, a_{j'}$ distincts, ce qui contredit le fait que les $a_j$ sont 
deux à deux premiers entre eux.
\end{proof}

\begin{thm}[des restes Chinois]
On suppose les $a_j$ deux à deux premiers entre eux, alors l'application :
\[ \varphi : x \in A \mapsto (\pi_1(x), \dots, \pi_r(x)) \in \prod_{j=1}^r A / (a_j) , \]
est un morphisme d'anneau surjectif de noyau $\ker{\varphi} = \left ( \prod a_j \right )$, 
et $\varphi$ induit un isomorphisme d'anneau :
\[ \veps : \pi(x) \in A / \left ( \prod a_j \right ) \mapsto (\pi_1(x), \dots, \pi_r(x)) \in \prod_{j=1}^r A / (a_j) , \]
Autrement dit, les anneaux $A / \prod a_j $ et $A / (a_1) \times \dots \times A / (a_r)$ sont isomorphes.
De plus, l'inverse de $\veps$ est donné par :
\[ \veps^{-1} : (\pi_1(x), \dots, \pi_r(x)) \in \prod_{j=1}^r A / (a_j) \mapsto 
	\widebar{\sum_{i=1}^r x_i u_i b_i } \in A / \left ( \prod a_j \right ), \]
où $(u_i)_{1 \leq i \leq r}$ est une suite d'éléments de $A$ telle que $\sum u_i b_i = 1$.
\end{thm}

\begin{proof}
En reprenant les notations introduites en début de développement, il est clair que l'application $\varphi$ est un morphisme d'anneaux.
Par ailleurs, un élément $x \in A$ est tel que $x \in \ker{\varphi}$ si, et seulement si, $\pi_j(x)=0$ pour tout $j$, ce qui équivaut à 
$x \in (a_j)$ pour tout $j$, soit $x \in (\ppcm(a_j))$. 
Comme les $a_j$ sont premiers entre eux, on a bien $\ker{\varphi} = \left ( \prod a_j \right )$. \\
Comme les $b_i$ sont premiers entre eux dans leur ensemble, le théorème de Bézout nous permet d'affirmer qu'il existe une suite  
$(u_i)_{1 \leq j \leq r}$ d'éléments de $A$ telle que $\sum u_i b_i = 1$. De plus, pour tour $1 \leq j \leq r$, si $i \neq j$ alors
$\pi_j(b_i) = \pi_j(0)$ (car $b_i$ est multiple de $a_j$), ce qui donne :
\[ 1 = \pi_j(1) = \pi_j \left( \sum_{i=1}^r u_i b_i \right) = \pi_j(u_j) \pi_j(b_j) , \]
autrement dit, $\pi_j(b_j)$ est inversible dans $A / (a_j)$ et son inverse est $\pi_j(u_j)$. \\
Ainsi, pour un $r$-uplet $(\pi_j(x_j))_j$ donné dans $\prod_j A / (a_j)$, en posant :
\[ x = \sum_{i=1}^r x_i u_i b_i , \qquad (*) \]
on a :
\[ \pi_j (x) = \pi_j(x_j) \pi_j(u_j) \pi_j(b_j) = \pi_j(x_j) , \]
ce qui démontre qe le morphisme $\varphi$ est surjectif. En quotientant $A$ par le noyau de $\varphi$, on obtient bien l'isomorphisme
$\veps$ annoncé. En considéant la classe d'équivalence (dans $A / \left ( \prod a_j \right ) )$ de $x$ défini par (*), 
on obtient bien la définition de la fonction inversé annoncée.
\end{proof}

\section{Application}

On veut résoudre le problème historique de Sun Zi, ce qui revient à résoudre le système modulaire :
\[ \left \{ \begin{array}{c} 
x \equiv 2 \pmod 3 \\ 
x \equiv 3 \pmod 5 \\ 
x \equiv 2 \pmod 7
\end{array} \right.  \]
Selon notre théorème, les anneaux $\Z/3\Z, \Z/5\Z$ et $\Z/7\Z$ sont isomorphes à l'anneau $\Z/105\Z$. 
De plus, résoudre le problème revient à trouver un inverse du triplet $(2,5,2)$.
Les trois \og produits partiels \fg{} (les $b_i$) sont 15, 21 et 35. On cherche un triplet $(u_1,u_2,u_3)$ tel que :
\[ 15x + 21y + 31z = 1 . \]
Comme $\gcd(35,21)=7$ et $\gcd(7,15)=1$, on cherche d'abord $a,b$ tels que $35a + 21b = 7$ puis $r,s$ tels que $7r + 15s = 1$.
On trouve :
\[ 2 \times 35 - 4 \times 21 + 15 = 1 \quad \text{soit } \left \{ \begin{array}{c} u_1 = 2 \\ u_2 = -4 \\ u_3 = 1 \end{array} \right.  \]
Il reste à calculer :
\[ x = 2 \times 2 \times 35 + 3 \times (-4) \times 21 + 2 \times 15 = -82 \]
Or $\widebar{-82} = 23$, donc le nombre d'objet est un élement de $\set{23 + 105 k ; k \in \N }$.

\end{document}